\exercisegroup

* 1. 设 E 是一个拓扑空间，并设 Σ 是 §2 练习 7 和 8 中定义的结构物种之一。取态射为这些练习中定义的态射，并取 α-映射为从 E 到 Σ-集的连续映射。证明 E 的（相对于前述定义的）泛映射问题一般没有解。*

* 2. 设 E 为一个域，Σ 为代数闭域结构的种。取态射为同态，α-映射为从 E 到代数闭域的同态。证明代数闭包 \(F_{E}\) E 的，以及 E 到 \(F_{E}\) 满足 \((AU_{I}^{\prime})\) ，但一般来说，E 的泛映射问题没有解。*

\begin{enumerate}
\def\labelenumi{\arabic{enumi}.}
\setcounter{enumi}{2}
\tightlist
\item
  让 \(\Sigma\) 设 为一个结构物种，设 \((\mathbf{A}_i)_{i \in I}\) 是一个集合族，并且对每个 \(i \in I\) 让 \(\mathcal{S}_i\) 是物种 \(\Sigma\) 在 \(\mathbf{A}_i\) 上的一个结构。设 E 为族 \((\mathbf{A}_i)_{i \in I}\) 的和集，其中每个 \(\mathbf{A}_i\) 被视为 E 的子集。假设我们给定了一个关于 \(\sigma\) -态射的概念，用于该物种 \(\Sigma\) ，并定义一个 \(\alpha\) 映射为一个映射 \(\varphi\) -集合 F 的 \(\Sigma\) -集 F，使得对于每个 \(i \in I\) ， \(\varphi\) 在 \(\mathbf{A}_i\) 上的限制是从 \(\mathbf{A}_i\) 到F中的一个态射。证明，如果存在一个解 \((\mathbf{F}_E, \varphi_E)\) 对于E的泛映射问题，则物种 \(\Sigma\) 在 \(\mathbf{F}_E\) 上的结构是关于该族 \((\mathbf{A}_i, \mathcal{S}_i, \varphi_i)_{i \in I}\) 的最终结构，其中 \(\varphi_i\) 表示 \(\varphi_E\) 在 \(\mathbf{A}_i\) 上的限制。
\end{enumerate}

此外，设 \(\mathbf{F}\) 设为一个集合，并且对于每个 \(\iota \in I\) 让 \(f_{\iota}\) 成为...的映射 \(\mathbf{A}_{\iota}\) 进入。 \(\mathbf{F}\) 。如果存在一个物种 \(\Sigma\) 在 \(\mathbf{F}\) 上、关于该族 \((\mathbf{A}_{\iota}, \mathcal{S}_{\iota}, f_{\iota})_{\iota \in I}\) 的最终结构，那么我们可以写 \(f_{\iota} = f \circ \varphi_{\iota}\) ，其中 \(f\) 是……的一个态射 \(\mathbf{F}_{\mathbf{E}}\) 进入。 \(\mathbf{F}\) ，并且其上的结构是 \(\mathbf{F}\) 在 \(f\) 下的直接像，该结构来自 \(\mathbf{F}_{\mathbf{E}}\) .

考虑以下情形的应用：

* (i) \(\Sigma\) 是一类代数结构，其态射为同态，且条件 \((\mathrm{CU}_{\mathbf{I}})\) 通过 \((\mathrm{CU}_{\mathbf{III}})\) 成立。这对于半群结构、群结构、模结构、代数结构等均成立。 \(\Sigma\) -集合 \(\mathbf{F}_{\mathbf{E}}\) 被称为 \(\mathbf{A}_{\mathbf{i}}\) 在群的情形下是直和，在模的情形下是直和，在代数的情形下是直积。

\begin{enumerate}
\def\labelenumi{(\roman{enumi})}
\setcounter{enumi}{1}
\tightlist
\item
  \(\Sigma\) 是拓扑群结构的种类或拓扑向量空间结构的种类。条件 \((\mathrm{CU}_{\mathrm{I}})\) 通过 \((\mathrm{CU}_{\mathrm{III}})\) 然后这些条件得到满足。对于局部凸拓扑向量空间的情形， \(\mathbf{F}_{\mathbf{E}}\) 被称为拓扑直和 \(\mathbf{A}_{i}\) . *
\end{enumerate}
